\documentclass[11pt]{article}
% ===========================================================================
%  Shared preamble for the AP Calculus AB exam-preparation worksheets.
%  Input right after \documentclass. Files beginning with "_" are not compiled
%  on their own. Navy and gold, to match the course's lesson pages.
% ===========================================================================
\usepackage[letterpaper,top=2.3cm,bottom=2.2cm,left=1.7cm,right=1.7cm,headheight=28pt,footskip=30pt]{geometry}
\usepackage{amsmath,amssymb}
\usepackage[T1]{fontenc}
\usepackage[utf8]{inputenc}
\usepackage{lmodern}
\usepackage{xcolor}
\usepackage[most]{tcolorbox}
\usepackage{enumitem}
\usepackage{fancyhdr}
\usepackage{array}
\usepackage{tikz}
\usepackage{pgfplots}
\pgfplotsset{compat=1.18}

\definecolor{navy}{HTML}{161E54}
\definecolor{ink}{HTML}{27357E}
\definecolor{gold}{HTML}{B98A00}
\definecolor{goldbg}{HTML}{FFFDF3}
\definecolor{inkbg}{HTML}{F6F7FD}
\definecolor{rule}{HTML}{C9CFF0}
\definecolor{okgreen}{HTML}{1B7A44}
\definecolor{okbg}{HTML}{F1F8F3}

\pagestyle{fancy}
\fancyhf{}
\fancyhead[L]{\footnotesize Name: \underline{\hspace{4cm}}}
\fancyhead[C]{\textbf{\color{navy}Integration Academy}\\[-2pt]{\scriptsize Calculus AB \textperiodcentered{} AP\textsuperscript{\textregistered} Exam Preparation}}
\fancyfoot[L]{\tiny\color{black!55}\parbox[t]{0.78\textwidth}{AP\textsuperscript{\textregistered} is a registered trademark of the College Board, which was not involved in the production of, and does not endorse, this material. \textcopyright{} Integration Academy.}}
\fancyfoot[R]{\scriptsize\thepage}
\renewcommand{\headrulewidth}{1.2pt}
\renewcommand{\headrule}{\hbox to\headwidth{\color{gold}\leaders\hrule height \headrulewidth\hfill}}
\renewcommand{\footrulewidth}{0pt}
\setlength{\parindent}{0pt}
\setlength{\parskip}{4pt}

% Title block: \wstitle{code}{title}{unit line}
\newcommand{\wstitle}[3]{\begin{center}{\small\bfseries\color{gold}WORKSHEET #1}\\[3pt]{\LARGE\bfseries\color{navy}#2}\\[3pt]{\small\color{black!60}#3}\end{center}\vspace{2pt}}
% Section heading
\newcommand{\sech}[1]{\par\vspace{10pt}{\large\bfseries\color{navy}#1}\par\nobreak\vspace{2pt}{\color{rule}\hrule height0.8pt}\vspace{7pt}}
% A level tag printed beside a group of questions
\newcommand{\level}[1]{\par\vspace{6pt}{\footnotesize\bfseries\color{gold}\MakeUppercase{#1}}\par\nobreak\vspace{2pt}}

\newtcolorbox{ideabox}{colframe=ink,colback=inkbg,coltitle=white,colbacktitle=ink,fonttitle=\bfseries\small,title={Key ideas},arc=2pt,boxrule=0.8pt,left=8pt,right=8pt}
\newtcolorbox{example}[1]{colframe=rule,colback=white,coltitle=navy,colbacktitle=inkbg,fonttitle=\bfseries\small,title={#1},arc=2pt,boxrule=0.8pt,borderline west={3pt}{0pt}{ink},left=10pt,right=8pt}
\newtcolorbox{frbox}[1]{colframe=gold,colback=goldbg,coltitle=white,colbacktitle=gold,fonttitle=\bfseries\small,title={#1},arc=2pt,boxrule=0.8pt,breakable,left=8pt,right=8pt}
\newtcolorbox{solbox}{colframe=okgreen,colback=okbg,coltitle=white,colbacktitle=okgreen,fonttitle=\bfseries\small,title={Worked solutions},arc=2pt,boxrule=0.8pt,breakable,left=8pt,right=8pt}
\newtcolorbox{scorebox}[1]{colframe=navy,colback=white,coltitle=white,colbacktitle=navy,fonttitle=\bfseries\small,title={#1},arc=2pt,boxrule=0.8pt,breakable,left=8pt,right=8pt}

\newcommand{\soln}{\par\smallskip\textbf{\color{ink}Solution.}\ }
\newcommand{\result}[1]{\par\smallskip\textbf{#1}}
% blank working space of a given height
\newcommand{\work}[1]{\par\vspace{#1}}
% a point value in the margin of a scoring guide
\newcommand{\pts}[1]{\hfill{\footnotesize\bfseries\color{gold}[#1]}}

% Coordinate grid for graph questions: \begin{apgraph}{xmin}{xmax}{ymin}{ymax} ... \end{apgraph}
\newenvironment{apgraph}[4]{\begin{center}\begin{tikzpicture}\begin{axis}[
  axis lines=middle, xmin=#1, xmax=#2, ymin=#3, ymax=#4,
  grid=both, grid style={black!12}, xtick distance=1, ytick distance=1,
  ticklabel style={font=\scriptsize}, xlabel={$x$}, ylabel={$y$},
  width=10cm, height=7.2cm, samples=80, clip=false,
  every axis plot/.append style={line width=1.1pt, color=ink}]}{\end{axis}\end{tikzpicture}\end{center}}
\newcommand{\closeddot}[2]{\addplot[only marks, mark=*, mark size=2.2pt, color=ink] coordinates {(#1,#2)};}
\newcommand{\opendot}[2]{\addplot[only marks, mark=*, mark size=2.2pt, color=ink, mark options={fill=white}] coordinates {(#1,#2)};}

% Multiple-choice options: four across, or stacked when the options are long.
\newcommand{\opts}[4]{\par\vspace{3pt}\begin{tabular}{@{}p{0.235\linewidth}p{0.235\linewidth}p{0.235\linewidth}p{0.235\linewidth}@{}}(A)\ #1 & (B)\ #2 & (C)\ #3 & (D)\ #4\end{tabular}\par\vspace{9pt}}
\newcommand{\optsv}[4]{\par\vspace{3pt}\begin{tabular}{@{}l@{\ \ }p{0.9\linewidth}@{}}(A) & #1\\(B) & #2\\(C) & #3\\(D) & #4\end{tabular}\par\vspace{9pt}}

\fancyhead[R]{\footnotesize\color{black!60}Worksheet 9.2}
\begin{document}

\wstitle{9.2}{A Phrase Bank for Justifications}{Unit 9: Exam Strategy and Free-Response Types}

\begin{ideabox}
\begin{itemize}[leftmargin=*,itemsep=2pt,topsep=2pt]
  \item \textbf{Extremum}: \textquotedblleft{}$f'$ changes from positive to negative at $x=c$, so $f$ has a relative maximum there.\textquotedblright{} Name the sign change.
  \item \textbf{Theorems}: state the condition, then the conclusion. IVT and EVT need continuity; MVT needs continuity on $[a,b]$ and differentiability on $(a,b)$.
  \item \textbf{Over or under}: a tangent line lies below a graph that is concave up; a left sum underestimates for an increasing function.
  \item \textbf{Speed}: increasing when $v$ and $a$ have the same sign. Give both values.
\end{itemize}
\end{ideabox}

\sech{Worked examples}
\begin{example}{Example 1 \textperiodcentered{} A sign change}
$f'(x)=(x-2)(x+1)^2$. Justify the location of every relative extremum of $f$.\soln $f'$ changes from negative to positive at $x=2$. At $x=-1$ it is negative on both sides.\result{A relative minimum at $x=2$ only; no extremum at $x=-1$ because $f'$ does not change sign.}
\end{example}

\begin{example}{Example 2 \textperiodcentered{} A theorem}
$f$ is continuous on $[1,5]$ with $f(1)=-3$ and $f(5)=4$. Must $f(c)=0$ for some $c$ in $(1,5)$?\soln $0$ lies between $f(1)$ and $f(5)$.\result{Yes: $f$ is continuous on $[1,5]$ and $-3<0<4$, so by the Intermediate Value Theorem such a $c$ exists.}
\end{example}

\begin{example}{Example 3 \textperiodcentered{} Over or under}
$f''(x)>0$ for all $x$. Is a tangent-line approximation of $f$ an overestimate or an underestimate?\soln The graph is concave up, so it lies above each tangent line.\result{An underestimate, because $f''>0$.}
\end{example}

\newpage
\sech{Practice}
\emph{No calculator on this page.}
\level{Core}
\begin{enumerate}[leftmargin=*,itemsep=2pt]
  \item $f'(3)=0$ and $f''(3)=-2$. What can you conclude, and why?\work{2.3cm}
  \item $g'(x)>0$ for all $x$ in $(0,4)$. Write a one-sentence justification that $g$ is increasing on $[0,4]$.\work{2.3cm}
  \item $h$ is differentiable on $[2,6]$ with $h(2)=1$ and $h(6)=13$. What does the Mean Value Theorem guarantee?\work{2.3cm}
  \item Is a left Riemann sum for a decreasing function an overestimate or an underestimate of the integral?\work{2.3cm}
\end{enumerate}

\level{Exam level}
\begin{enumerate}[leftmargin=*,itemsep=2pt,start=5]
  \item $f$ is differentiable for all $x$, with $f(0)=5$, $f(2)=1$ and $f(5)=7$. Explain why $f'(c)=0$ for some $c$ in $(0,5)$.\work{2.8cm}
  \item $f''(x)=(x-1)^2(x-4)$. Find the $x$-coordinate of every point of inflection of $f$, with justification.\work{2.8cm}
  \item $f''(x)>0$ on $[a,b]$. Is a trapezoidal sum for $\displaystyle\int_{a}^{b}f(x)\,dx$ an overestimate or an underestimate? Give a reason.\work{2.8cm}
  \item A particle has $v(3)=-4$ and $a(3)=2$. Is its speed increasing or decreasing at $t=3$? Give a reason.\work{2.8cm}
\end{enumerate}

\begin{frbox}{Question 9 \textperiodcentered{} Free-response practice}
The function $f$ is twice differentiable. Selected values of $f$ and $f'$ are given in the table.\begin{center}\renewcommand{\arraystretch}{1.25}\begin{tabular}{|c|c|c|c|c|}\hline $x$ & $1$ & $3$ & $5$ & $8$ \\ \hline $f(x)$ & $2$ & $6$ & $5$ & $11$ \\ \hline $f'(x)$ & $3$ & $1$ & $-2$ & $4$ \\ \hline\end{tabular}\end{center}
\begin{enumerate}[label=(\alph*),leftmargin=*,itemsep=2pt]
  \item Must there be a value $c$, $3<c<5$, with $f'(c)=-\dfrac12$? Justify your answer.\work{2.6cm}
  \item Must there be a value $c$, $5<c<8$, with $f'(c)=0$? Justify your answer.\work{2.6cm}
  \item Use the tangent line at $x=3$ to approximate $f(3.5)$. Given that $f''(x)<0$ for $3<x<5$, is the approximation an overestimate or an underestimate? Give a reason.\work{2.6cm}
\end{enumerate}
\end{frbox}

\newpage
\begin{solbox}
\begin{enumerate}[leftmargin=*,itemsep=4pt]
  \item $f$ has a relative maximum at $x=3$, by the Second Derivative Test: $f'(3)=0$ and $f''(3)<0$.
  \item $g$ is increasing on $[0,4]$ because $g'(x)>0$ on $(0,4)$.
  \item $h$ is differentiable, hence continuous, so there is a $c$ in $(2,6)$ with $h'(c)=\dfrac{13-1}{6-2}=3$.
  \item An overestimate: each rectangle uses the left, largest, value.
  \item $f$ is continuous, and $5$ lies between $f(2)=1$ and $f(5)=7$, so by the IVT $f(b)=5$ for some $b$ in $(2,5)$. Then $f(0)=f(b)$, and the MVT on $[0,b]$ gives $f'(c)=0$.
  \item $x=4$ only: $f''$ changes from negative to positive there. At $x=1$, $f''$ is negative on both sides, so there is no change of concavity.
  \item An overestimate: the graph is concave up, so each chord lies above the graph.
  \item Decreasing, because $v(3)<0$ and $a(3)>0$ have opposite signs.
  \item \begin{enumerate}[label=(\alph*),leftmargin=*,itemsep=3pt]
    \item Yes. $f$ is differentiable, hence continuous, on $[3,5]$, so by the MVT there is a $c$ with $f'(c)=\dfrac{5-6}{5-3}=-\dfrac12$.
    \item Yes. $f$ is twice differentiable, so $f'$ is continuous on $[5,8]$, and $f'(5)=-2<0<4=f'(8)$. By the IVT, $f'(c)=0$ for some $c$.
    \item $f(3.5)\approx6+1(0.5)=6.5$. An overestimate: $f''<0$, so the graph is concave down and lies below its tangent line.
  \end{enumerate}
\end{enumerate}
\end{solbox}

\end{document}
